\section{Back-Pressure：自动验证机制}

\subsection{核心洞察}

HumanLayer 的核心洞察：\textbf{使用 coding agent 成功解决问题的概率，与 agent 验证自身工作的能力强相关。}

他们在构建测试和其他 back-pressure 机制上投入了大量时间，这仍然是投资回报率最高的工作之一。

\subsection{验证机制的类型}

HumanLayer 代码库中的验证机制包括：

\begin{table}[H]
\centering
\caption{Back-Pressure 验证机制}
\begin{tabular}{@{}p{4cm}p{9cm}@{}}
\toprule
\textbf{机制} & \textbf{说明} \\
\midrule
类型检查和构建步骤 & 最好使用强类型语言 \\
单元测试/集成测试 & 验证功能正确性 \\
代码覆盖率报告 & 有一个 Stop hook 在覆盖率下降时提示 agent 增加覆盖率 \\
UI 交互和测试集成 & Playwright、agent-browser 等 \\
\bottomrule
\end{tabular}
\end{table}

\subsection{验证机制的 Context Efficiency}

\begin{warningbox}{验证输出必须精简---血泪教训}
HumanLayer 早期让 agent 在每次改动后运行完整的测试套件，结果 4000 行通过的测试输出淹没了 context window。Agent 随后失去了对实际任务的跟踪，开始对它刚刚读取的测试文件产生幻觉。

现在的做法：\textbf{吞掉成功输出，只呈现错误}。构建也是一样---\textbf{成功是静默的，只有失败才产生详细输出}。

验证机制必须是 context-efficient 的，否则验证本身就成了干扰源。
\end{warningbox}

HumanLayer 在 CLAUDE.md 文件中给 Claude 提供了关于如何使用所有这些验证机制的简洁指令，有些甚至被打包进 skills 中以实现渐进式披露。

\subsection{本章小结}

Back-pressure 是提高 agent 任务成功率的最高杠杆点之一。关键是让 agent 能自我验证，但验证输出必须精简：成功静默，失败详报。否则验证本身会污染 context window，反而降低 agent 的表现。


